\section{Statistical Procedures}
\label{app:stats}

This appendix defines every statistic in the paper. All tests are two-sided, and all code uses only the Python standard library, with random seed 0.

\paragraph{Page effect.}
For a question $q$ and an outcome $o\in\{\text{correct},\text{wrong},\text{decline}\}$, let $y_q^{o,\text{page}}=1$ if the model's answer has outcome $o$ when it reads the page, and $0$ otherwise. Let $y_q^{o,\text{closed}}$ be the same without the page. The page effect of $q$ is
\begin{equation}
d_q^{o}=y_q^{o,\text{page}}-y_q^{o,\text{closed}}\in\{-1,0,1\}.
\end{equation}
A page effect of $+1$ for declines means the page turned an answer into a decline; $-1$ means the page turned a decline into an answer. Our main measure (Eq.~\ref{M-eq:did}) is the mean page effect on family-only questions $F$ minus that on silent questions $S$:
\begin{equation}
\Delta^{o}=\bar d^{\,o}_F-\bar d^{\,o}_S,\qquad \bar d^{\,o}_G=\frac{1}{|G|}\sum_{q\in G}d_q^{o}.
\end{equation}
This is a difference in differences. The first difference (page minus closed book) removes how hard each fact is on its own. The second (family-only minus silent) asks what the family name adds beyond a page that says nothing.

\paragraph{Pooling models.}
For the pooled 7--8B result, we first average each question's page effect over Qwen2.5-7B, OLMo~2 7B and Granite~3.1 8B, and then compute $\Delta^o$ on these averages. Each of the 481 questions thus counts once, however many models answer it.

\paragraph{Confidence intervals.}
We draw 5,000 bootstrap samples. Each sample redraws the questions of each group with replacement, keeping the group sizes, and recomputes $\Delta^o$. The 95\% CI is the 2.5th to 97.5th percentile of these values. If the interval excludes zero, we call the effect reliable.

\paragraph{$p$-values.}
We test $\Delta^o=0$ with a Welch-type $z$-test, which allows the two groups to have different variances:
\begin{equation}
z=\frac{\bar d^{\,o}_F-\bar d^{\,o}_S}{\sqrt{s_F^2/|F|+s_S^2/|S|}},
\end{equation}
where $s_G^2$ is the sample variance of $d_q^o$ in group $G$. The $p$-value is $2\,(1-\Phi(|z|))$, with $\Phi$ the standard normal distribution function. A small $p$ means a gap this large is unlikely if the family name had no effect. Each comparison covers three outcomes, so the Bonferroni-corrected $p$-value is three times the raw one.

\paragraph{Odds ratios.}
For the closed-book check in \cref{app:data}, each relation $r$ gives a $2\times2$ table: $a_r$ and $b_r$ are correct and wrong answers in the first group, $c_r$ and $e_r$ in the second, and $n_r$ is the total. The Mantel--Haenszel odds ratio pools the relations:
\begin{equation}
\mathrm{OR}_{\mathrm{MH}}=\frac{\sum_r a_r e_r/n_r}{\sum_r b_r c_r/n_r}.
\end{equation}
$\mathrm{OR}=1$ means the two groups are answered correctly equally often. $\mathrm{OR}=0.27$ means the odds of a correct answer in the first group are about a quarter of those in the second. Pooling within relations keeps the result from being driven by one relation that happens to be easier.

\paragraph{AUROC.}
To judge confidence filtering, we take the answers a model gives (not its declines) and ask how well the probability of the chosen option ranks correct answers above wrong ones. AUROC is the chance that a random correct answer has higher confidence than a random wrong answer. An AUROC of 1.0 means a threshold can separate them perfectly; 0.5 means confidence carries no signal.

\paragraph{McNemar's test.}
In recovery, two choosers answer the same 410 facts. Let $b$ be the number of facts the first gets right and the second gets wrong, and $c$ the reverse. McNemar's test uses only these disagreements, with a continuity correction:
\begin{equation}
\chi^2=\frac{(|b-c|-1)^2}{b+c},
\end{equation}
compared with a $\chi^2$ distribution with one degree of freedom. It asks whether one chooser wins the disagreements more often than chance would allow.
